% Clase 26 --- Principio de Huygens (§4).
% (a) el ejemplo del docente: la cubeta golpeada con un objeto CUADRADO; los
%     frentes se van redondeando hasta parecer circulares.
% (b) el limite del metodo: partiendo de un frente plano la envolvente es otro
%     plano, y el principio no aporta nada.
\begin{center}
\begin{tikzpicture}[scale=0.92, transform shape]

% ---------------- (a) el cuadrado que se redondea ----------------
\begin{scope}
  % frente actual: cuadrado
  \draw[figline, line width=1pt] (-0.85,-0.85) rectangle (0.85,0.85);

  % ondas secundarias: un circulito centrado en cada punto del frente
  \foreach \pp in {(-0.85,-0.85),(-0.85,-0.28),(-0.85,0.28),(-0.85,0.85),
                   (-0.28,0.85),(0.28,0.85),(0.85,0.85),(0.85,0.28),
                   (0.85,-0.28),(0.85,-0.85),(0.28,-0.85),(-0.28,-0.85)}{
    \draw[figaux] \pp circle (0.55);}

  % envolvente: nuevo frente
  \draw[figamber, line width=1pt, rounded corners=0.55cm]
    (-1.4,-1.4) rectangle (1.4,1.4);
  % un paso mas
  \draw[figamber!55, line width=1pt, rounded corners=1.35cm]
    (-2.15,-2.15) rectangle (2.15,2.15);

  \node[figlbl, figamber, anchor=south west] at (0.95,1.45) {envolvente};

  \node[figlbl, anchor=north, align=center] at (0,-2.55)
    {\textbf{(a)} el frente cuadrado se redondea\\
     paso a paso (radio $=\lambda$ por período)};
\end{scope}

% ---------------- (b) el caso en que no aporta nada ----------------
\begin{scope}[xshift=7.4cm]
  \draw[figline, line width=1pt] (-1.7,-0.6) -- (1.7,-0.6);
  \foreach \xx in {-1.7,-1.13,...,1.75}{
    \draw[figaux] (\xx,-0.6) circle (0.55);}
  \draw[figamber, line width=1pt] (-1.7,-0.05) -- (1.7,-0.05);

  \draw[figvecaux] (0,-0.6) -- (0,-0.05);
  \node[figlblaux, anchor=west] at (0.08,-0.32) {$\lambda$};

  \node[figlbl, anchor=north, align=center] at (0,-2.55)
    {\textbf{(b)} si el frente ya es plano,\\
     la envolvente es otro plano};
\end{scope}

\end{tikzpicture}\\[0.5em]
{\small\itshape Formulado en 1678 ---mucho antes de que se supiera que la luz es
un campo electromagnético--- el principio vale para \textbf{cualquier} onda y es
una receta gráfica, no una fórmula: cada punto del frente actual se comporta como
una fuente puntual de ondas esféricas, y el frente siguiente es la
\textbf{envolvente} de todos esos circulitos. Como todas las ondas secundarias
tienen la misma $\lambda$ y el mismo período, \textbf{todos los circulitos son
iguales}; el radio depende de cuánto se avanza por paso ($\lambda$ si es un
período entero, $\lambda/10$ si es una décima parte). Su valor está en los frentes
\textbf{curvos}, donde permite predecir la deformación con regla y compás: en el
panel (b) se ve que para un frente plano el método sólo confirma que la onda sigue
siendo plana. Tiene una falla conocida: tomado al pie de la letra también predice
un frente que viaja \textbf{hacia atrás}. La salida es que las ondas secundarias
no son esféricas uniformes sino que tienen una dirección preferencial ---para
saber hacia dónde va la onda, y no sólo dónde está, hay que conocer cuánta
intensidad manda hacia cada lado.}
\end{center}
